\documentclass[11pt, a4paper]{article}

% Paquetes esenciales
\usepackage[utf8]{inputenc}
\usepackage[spanish, es-tabla]{babel} % Idioma español y renombrar "Cuadro" a "Tabla"
\usepackage{amsmath, amsfonts, amssymb} % Matemáticas
\usepackage{graphicx} % Gráficos
\usepackage{geometry} % Márgenes
\usepackage{booktabs} % Tablas profesionales
\usepackage{caption}
\usepackage{float} % Posicionamiento de figuras
\usepackage[colorlinks=true, linkcolor=blue, citecolor=blue, urlcolor=blue]{hyperref} % Enlaces
\usepackage{fancyhdr} % Encabezados
\usepackage{siunitx} % Unidades del SI
\sisetup{output-decimal-marker = {.}}

% Configuración de márgenes
\geometry{top=2.5cm, bottom=2.5cm, left=2.5cm, right=2.5cm}

% Configuración del encabezado (estilo profesional P4)
\pagestyle{fancy}
\fancyhf{}
\lhead{Termodinámica - Grado en Física}
\rhead{Práctica 7: Método de Clement-Desormes}
\cfoot{\thepage}

% Título y datos del documento
\title{\vspace{-2cm}\textbf{Práctica 7: Determinación del coeficiente adiabático de gases \\ Método de Clement-Desormes}}
\author{
    \textbf{Andrés Vinuesa Espinosa} y \textbf{José María Martínez Herrada} \\
    Grupo B2.4 \\
    Universidad de Granada
}
\date{Realización: 08/04/2026 \quad | \quad Entrega: 22/04/2026}

\begin{document}

\maketitle

\hrule
\vspace{0.5cm}

\begin{abstract}
\noindent 
En este informe se determina experimentalmente el coeficiente adiabático ($\gamma$) del aire y del dióxido de carbono ($CO_2$) utilizando el método de Clement-Desormes. Se analiza la dependencia de $\gamma$ con respecto a la presión inicial aplicada en el sistema y se comparan los resultados empíricos con los valores teóricos deducidos a partir de los grados de libertad de moléculas diatómicas y lineales poliatómicas. Finalmente, se discuten las contribuciones rotacionales y vibracionales en el caso del $CO_2$ \cite{guion}.
\end{abstract}

\vspace{0.5cm}
\hrule

\vspace{1cm}
\tableofcontents
\newpage

% ---------------------------------------------------------
% 1. RESULTADOS
% ---------------------------------------------------------
\section{Resultados}

\subsection{Coeficiente adiabático del aire}
En esta sección calculamos el valor del coeficiente adiabático $\gamma$ para el aire a diferentes presiones iniciales. Utilizando la aproximación del método de Clement-Desormes para variaciones de presión pequeñas frente a la atmosférica, la relación teórica viene dada por la ecuación:

\begin{equation}
    \gamma \approx \frac{h_1}{h_1 - h_2}
    \label{eq:gamma_aprox}
\end{equation}

En la Tabla \ref{tab:resultados_aire} (Anexo \ref{anexo:tablas}) se recogen los valores obtenidos.

\begin{figure}[H]
    \centering
    \includegraphics[width=0.7\textwidth]{grafica_aire.png}
    \caption{Gráfico de $\gamma$ frente a $h_1$ para el aire.}
    \label{fig:gamma_aire}
\end{figure}

Al observar los resultados, vemos que los valores de $\gamma$ experimentan cierta dispersión dependiendo de la diferencia de altura inicial, oscilando entre $1.31$ y $1.46$. Teóricamente, el coeficiente adiabático es una propiedad intrínseca del gas dependiente de la temperatura, por lo que debería mantenerse constante frente a variaciones de presión inicial en un gas ideal. Las diferencias observadas de los $\gamma$ observados se podrían atribuir a errores instrumentales en las lecturas más pequeñas (como en $h_1=6.4$ cm) y al error experimental al no hacer la expansión adiábatica con unos instrumentos y sistemas de la máxima calidad.

Suponiendo que el aire se comporta como un gas ideal diatómico puro (compuesto principalmente de $N_2$ y $O_2$), sus grados de libertad son $f=5$ (3 traslacionales y 2 rotacionales a temperatura ambiente)\cite{young}. Por tanto, el valor teórico esperado es:

\begin{equation}
    \gamma_{teo} = \frac{C_p}{C_v} = \frac{f+2}{f} = \frac{7}{5} = 1.40
\end{equation}

Calculando la media de nuestros valores experimentales, obtenemos un coeficiente de $\gamma_{exp} = 1.41 \pm 0.35$. Al compararlo con el valor teórico y bibliográfico ($\gamma_{bib} = 1.40 - 1.41$) \cite{young}, comprobamos que las mediciones concuerdan bastante bien con las predicciones, presentando un error menor al $1\%$. Esto corrobora el modelo teórico de gas diatómico para el aire a temperatura ambiente.

\subsection{Coeficiente adiabático del $CO_2$}
A continuación, repetimos el procedimiento habiendo purgado y llenado el sistema con dióxido de carbono ($CO_2$). Los valores exactos calculados y medidos se exponen en la Tabla \ref{tab:resultados_co2}.

\begin{figure}[H]
    \centering
    \includegraphics[width=0.7\textwidth]{grafica_co2.png}
    \caption{Gráfico de $\gamma$ frente a $h_1$ para el $CO_2$.}
    \label{fig:gamma_co2}
\end{figure}

Al igual que con el aire, $\gamma$ debería ser independiente de $h_1$. Sin embargo, encontramos valores que van desde $1.27$ hasta $1.46$. Estas variaciones responden a la mayor dificultad experimental para estabilizar térmica y mecánicamente un gas más denso como el $CO_2$, así como a posibles mezclas residuales de aire si el purgado del sistema no fue perfectamente hermético. 

Considerando el $CO_2$ como una molécula lineal poliatómica, posee 3 grados de libertad traslacionales y 2 rotacionales. Si asumiésemos la ausencia total de contribución vibracional a temperatura ambiente ($f=5$), su valor teórico sería idéntico al del aire ($\gamma = 1.40$). La bibliografía, no obstante, sitúa el $\gamma$ del $CO_2$ en $1.294$ \cite{young}. Aunque el $CO_2$ es una molécula lineal similar a las del aire, el valor que nos dice en la bibliografía es de $\gamma \approx 1.3$ refleja la contribución de los modos vibracionales que aumentan los grados de libertad.

Nuestra media experimental se sitúa en $\gamma_{exp} = 1.37 \pm 0.12$. Aunque superior al tabulado, es menor que $1.40$. El hecho de que nuestros resultados experimentales se aproximen a $1.4$ sugiere que, bajo las presiones y temperaturas del laboratorio, el gas se ha comportado parecido al aire, omitiendo la contribución vibracional que predice la teoría.

% ---------------------------------------------------------
% 2. CONCLUSIONES
% ---------------------------------------------------------
\section{Conclusiones}

En esta práctica se ha determinado el coeficiente adiabático de dos gases de distinta naturaleza molecular (aire y $CO_2$) mediante el método de Clement-Desormes, analizando cómo influyen los grados de libertad en la capacidad calorífica de los fluidos.

Para el caso del aire, el resultado experimental nos da un valor medio de $\gamma = 1.41 \pm 0.35$, el cual se ajusta perfectamente a las predicciones teóricas y bibliográficas ($\gamma \approx 1.40$) con un error inferior al 1\,\%. Por lo tanto, los datos tomados en la práctica corroboran el modelo de gas ideal diatómico para el aire a temperatura ambiente, demostrando que la energía térmica se distribuye de forma exclusiva entre sus 5 grados de libertad (traslacionales y rotacionales).

Por su parte, el análisis para el $CO_2$ proporcionó un promedio de $\gamma = 1.37 \pm 0.12$. Aunque la bibliografía predice un valor cercano a $1.294$ debido a la activación de sus modos vibracionales, nuestro resultado se aproxima de forma notable al $1.40$. Esto sugiere que, bajo las presiones y temperaturas específicas de nuestro laboratorio, la predicción dela teoría ha quedado omitida, haciendo que el gas poliatómico se comporte de manera muy similar al aire.

Finalmente, se concluye que la dispersión observada en los datos a distintas presiones iniciales refleja las limitaciones propias del diseño experimental. Las desviaciones encontradas se atribuyen a la dificultad práctica de aislar térmicamente el sistema para lograr una expansión estrictamente adiabática con el equipo disponible, así como al peso del error instrumental en las lecturas más pequeñas ($h_1=6.4$ cm). En el caso específico del $CO_2$, el margen de error se vio agravado por la muy probable mezcla de aire residual derivada de un purgado no perfectamente hermético.
\newpage

\section{Anexos}

\subsection{Tablas de datos}
\label{anexo:tablas}

\begin{table}[H]
\centering
\caption{Alturas promedio del líquido manométrico para el \textbf{aire} y cálculo exacto de $\gamma$.}
\vspace{0.2cm}
\begin{tabular}{cccc}
\toprule
\textbf{Medida} & \textbf{$h_1$ (cm)} & \textbf{$h_2$ (cm)} & \textbf{$\gamma_{aire}$ exacto} \\
\midrule
1 & $24.60\pm 0.10$  & $7.20\pm 0.10$  & $1.41 \pm 0.20$ \\ 
2 & $12.50\pm 0.10$ & $3.80\pm 0.10$ & $1.44 \pm 0.42$ \\
3 & $35.80\pm 0.10$ & $11.30\pm 0.10$ & $1.46 \pm 0.15$ \\
4 & $6.40\pm 0.10$ & $1.50\pm 0.10$ & $1.31 \pm 0.61$ \\
\bottomrule
\end{tabular}
\label{tab:resultados_aire}
\end{table}

\begin{table}[H]
\centering
\caption{Alturas promedio del líquido manométrico para el \textbf{$CO_2$} y cálculo exacto de $\gamma$.}
\vspace{0.2cm}
\begin{tabular}{cccc}
\toprule
\textbf{Medida} & \textbf{$h_1$ (cm)} & \textbf{$h_2$ (cm)} & \textbf{$\gamma_{CO_2}$ exacto} \\
\midrule
1 & $32.10 \pm 0.10$ & $9.10 \pm 0.10$ & $1.394 \pm 0.094$ \\ 
2 & $14.40 \pm 0.10$ & $3.10 \pm 0.10$ & $1.27 \pm 0.16$ \\
3 & $24.50 \pm 0.10$ & $6.50 \pm 0.10$ & $1.36 \pm 0.12$ \\
4 & $28.60 \pm 0.10$ & $9.00 \pm 0.10$ & $1.46 \pm 0.12$ \\
\bottomrule
\end{tabular}
\label{tab:resultados_co2}
\end{table}


\subsection{Incertidumbres}
\label{anexo:incertidumbres}
A la hora de realizar las mediciones, nos enfrentamos a dos tipos de incertidumbres: directas e indirectas.

\paragraph{Medidas Directas (Incertidumbre Tipo B):} 

Las únicas magnitudes medidas de forma directa en el laboratorio han sido las diferencias de altura en el tubo manométrico ($h_1$ y $h_2$). Su error instrumental se asocia a la lectura en escala graduada. En este caso siendo 1mm el error del instrumento.

\paragraph{Medidas Indirectas (Propagación Cuadrática de Errores):}
El coeficiente adiabático $\gamma$ es una medida indirecta. Para calcular su incertidumbre $\Delta\gamma$, aplicamos la ley general de propagación cuadrática de errores para variables independientes:
\begin{equation}
    \Delta \gamma = \sqrt{ \left( \frac{\partial \gamma}{\partial h_1} \Delta h_1 \right)^2 + \left( \frac{\partial \gamma}{\partial h_2} \Delta h_2 \right)^2 }
\end{equation}

Utilizando la ecuación aproximada de Taylor $\gamma \approx \frac{h_1}{h_1 - h_2}$, el desarrollo analítico explícito de las derivadas parciales resulta en:

\begin{equation}
    \frac{\partial \gamma}{\partial h_1} = \frac{1 \cdot (h_1 - h_2) - h_1 \cdot 1}{(h_1 - h_2)^2} = \frac{-h_2}{(h_1 - h_2)^2}
\end{equation}

\begin{equation}
    \frac{\partial \gamma}{\partial h_2} = - \frac{h_1}{(h_1 - h_2)^2} \cdot (-1) = \frac{h_1}{(h_1 - h_2)^2}
\end{equation}

Sustituyendo las derivadas en la ecuación de propagación, obtenemos la expresión analítica final aproximada que rige la dinámica de los errores presentados en las tablas:
\begin{equation}
    \Delta \gamma \approx \frac{1}{(h_1 - h_2)^2} \sqrt{ (-h_2 \Delta h_1)^2 + (h_1 \Delta h_2)^2 }
\end{equation}



% ---------------------------------------------------------
% REFERENCIAS 
% ---------------------------------------------------------
\begin{thebibliography}{9}
\bibitem{guion} 
Departamento de Termódinamica. \textit{Guion Práctica 7: Determinación del coeficiente adiabático de gases: método de Clement-Desormes}. Universidad de Granada, Curso 2025-2026.

\bibitem{tipler} 
Tipler, P. A., \& Mosca, G. (2010). \textit{Física para la ciencia y la tecnología. Vol. 1: Mecánica, oscilaciones y ondas, termodinámica} (6ª ed.). Reverté.

\bibitem{young}
Young, H. D., \& Freedman, R. A. (2013). \textit{Física Universitaria con Física Moderna} (13ª ed., Vol. 1). Pearson Educación.

\end{thebibliography}

\end{document}